\documentclass[11pt]{article}
\usepackage{preamble}
\setlength{\parskip}{4pt}

\begin{document}

\daybanner{AP Statistics \textperiodcentered\ Unit 1 \textperiodcentered\ Topic 1.1 \textperiodcentered\ B: 2026-09-08 \textperiodcentered\ E: 2026-09-09 \textperiodcentered\ TEACHER KEY}{What Can These Data Support?}

\sectionbanner{CED TETHER}

{\small
\begin{itemize}[leftmargin=1.5em,itemsep=2pt,topsep=2pt]
  \item \textbf{VAR-1.A} Identify questions to be answered, based on variation in one-variable data. \textbf{[Skill 1.A]}
  \item \textbf{VAR-1.A.1} Numbers may convey meaningful information, when placed in context.
\end{itemize}
}
\textbf{Essential question:} How do the data collected limit what we can honestly claim?\par
\textbf{DOK-3 skill:} 4.A \textperiodcentered\ \textbf{FRQ pattern:} {\small\texttt{claim-\allowbreak{}vs-\allowbreak{}data-\allowbreak{}support-\allowbreak{}bound-\allowbreak{}the-\allowbreak{}conclusion}}\par
\textbf{DOK rationale:} Part (c) coordinates a paired-data pattern with selection and collection methods, bounds a descriptive conclusion, and matches new evidence to a missing scope or causal link.\par

\frameworkphaseheader{Do Now (first take) $\rightarrow$ Explore (video + follow-along) $\rightarrow$ Exit (finish + turn in)}{1 $\rightarrow$ 3}{45}{%
\begin{itemize}[leftmargin=1.5em,itemsep=2pt,topsep=2pt]
  \item Board slide up at the bell. Read both newspaper claims once without endorsing either.
  \item Begin the video follow-along at minute 5; leave first takes ungraded.
  \item Return to the board slide at minute 33. Circulate and ask students to separate pattern, scope, and cause.
  \item Collect at the door. Score part (c) only, E/P/I.
\end{itemize}
}{%
\begin{itemize}[leftmargin=1.5em,itemsep=2pt,topsep=2pt]
  \item Commit to one claim, neither claim, or both in a one-sentence first take.
  \item Complete the video follow-along about questions, variation, and context.
  \item Finish (a)--(c), revise the first take in the exit line, and turn in the sheet.
\end{itemize}
}{%
\begin{itemize}[leftmargin=1.5em,itemsep=2pt,topsep=2pt]
  \item Which paired values make the words 'tended to' reasonable even with Leo in the table?
  \item Which words in Claim 2 change the scope, and which word makes it causal?
  \item What does 'surveyed' tell you that 'randomly assigned' would not?
  \item Would a random school sample establish cause? Would random assignment in this one class represent the school?
\end{itemize}
}{Solo teacher. Circulate ~5 pairs. Ask students to point to the table or collection description before accepting a claim; do not supply the words 'scope' or 'cause' first.}

\begin{callout}{\IconWarn\ WATCH FOR}{calloutred}
\begin{itemize}[leftmargin=1.5em,itemsep=2pt,topsep=2pt]
  \item Using Leo to erase the overall tendency instead of treating him as one exception.
  \item Saying only '12 is too small' without explaining who the 12 students represent.
  \item Treating a random sample as proof of cause, or random assignment in one class as proof about the whole school.
\end{itemize}
\end{callout}

\sectionbanner{SCORING PART (c) --- E / P / I}

\textbf{Expected elements}
\begin{itemize}[leftmargin=1.5em,itemsep=2pt,topsep=2pt]
  \item \textbf{supports-claim-1-from-pattern} (required) --- Supports Claim 1 as a description and cites the upward tendency, allowing the exception
  \item \textbf{rejects-claim-2-on-scope} (required) --- Rejects school-wide scope because one class is not a representative school sample
  \item \textbf{rejects-claim-2-on-cause} (required) --- Rejects causation because sleep was observed, not assigned, with no comparison
  \item \textbf{matches-new-data-to-limit} (required) --- Names specific additional data and says whether it addresses scope, causation, or both
\end{itemize}

{\small\renewcommand{\arraystretch}{1.25}
\noindent\begin{tabular}{|p{0.5in}|p{5.9in}|}
\hline
\textbf{E} & Supports Claim 1 from the table, rejects Claim 2 with both scope and causal limits, and matches a data plan to a limit. \\ \hline
\textbf{P} & Makes both decisions but explains only one limit, proposes data without naming its purpose, or cites only the exception. \\ \hline
\textbf{I} & Supports Claim 2, rejects Claim 1 solely because of Leo, or gives no data or collection evidence. \\ \hline
\end{tabular}}

\textbf{Common mistakes}
\begin{itemize}[leftmargin=1.5em,itemsep=2pt,topsep=2pt]
  \item Rejecting Claim 1 because of Leo although ``tended'' allows an exception
  \item Saying only ``12 is too small'' instead of explaining whom one class represents
  \item Treating an observed association as proof that sleep caused the scores
\end{itemize}

\clearpage
\focusbanner{TODAY'S PROBLEM --- annotated}

A newspaper asked the 12 students in one AP Statistics class about sleep before a 10-point quiz; no sleep amount was assigned. The table shows the pairs.\par\smallskip \textbf{Claim 1:} ``In this class, students who slept more tended to score higher.''\par \textbf{Claim 2:} ``Sleeping more raises quiz scores for all students at our school.''\par
\begin{center}
\textbf{Sleep and quiz scores}\par\smallskip
\begin{tabular}{|l|c|c|}
\hline
\textbf{Student} & \textbf{Sleep (h)} & \textbf{Quiz (/10)} \\ \hline
\textbf{Aiden} & 5.0 & 4 \\ \hline
\textbf{Bea} & 7.0 & 7 \\ \hline
\textbf{Cruz} & 8.5 & 10 \\ \hline
\textbf{Dev} & 6.5 & 7 \\ \hline
\textbf{Elise} & 8.0 & 8 \\ \hline
\textbf{Fiona} & 7.5 & 8 \\ \hline
\textbf{Gabe} & 5.5 & 5 \\ \hline
\textbf{Hana} & 9.0 & 10 \\ \hline
\textbf{Imani} & 7.0 & 7 \\ \hline
\textbf{Jalen} & 6.0 & 6 \\ \hline
\textbf{Kira} & 8.0 & 9 \\ \hline
\textbf{Leo} & 4.5 & 9 \\ \hline
\end{tabular}
\end{center}
\begin{firsttakebox}{FIRST TAKE (5 min)}
Which claim, if either, would you print? Name one detail behind your choice.\par
\answer{Not graded. Expect Claim 2; the revision should become more bounded.}
\end{firsttakebox}

\textit{Rules callout ``CLAIMS NEED CONTEXT AND EVIDENCE (keep on the page)'' is printed on the student sheet.}\par\medskip

\dokbadge{1}~\textbf{(a)} State the statistical question and identify both variables, with units.\par
\answer{Question: For this class, is quiz score related to sleep the night before? Variables: sleep (hours) and quiz score (points out of 10).}

\dokbadge{2}~\textbf{(b)} Describe the relationship for these students. Cite values for the pattern and its exception.\par
\answer{Except for Leo, more sleep goes with higher scores: 5--6 hours pairs with 4--6 points, while 8--9 hours pairs with 8--10. Leo slept 4.5 hours and scored 9.}

\dokbadge{3}~\textbf{(c)} Which claim is supported? Cite table evidence for Claim 1. Explain how one class and no assigned comparison limit Claim 2. Name one kind of additional data and the limit it would address.\par
\answer{Claim 1 is supported for these 12: most pairs rise, and ``tended'' allows Leo. Claim 2 is not: one class does not represent the school, and observation with no assigned comparison cannot establish cause. A representative school sample addresses scope only; a randomized comparison of a safe sleep-support intervention addresses cause for its participants only. Accept either with its limit named.}

\begin{summaryexitbox}{EXIT}
One thing the video changed about my first take: \hrulefill
\end{summaryexitbox}

\end{document}
